\section{Related Work} \label{ref}

\subsection{Segmentation of remote sensing imagery} \label{subsec:seg-rsd}
Early semantic segmentation models were largely based on convolutional neural networks (CNNs) \cite{7789580}. A key milestone was the fully convolutional network (FCN) introduced in \cite{FCN}, which adopts an encoder–decoder architecture composed entirely of convolutional layers, without any fully connected layers. Building on this foundation, \cite{rs13234902} proposed an improved FCN framework for remote sensing image segmentation, incorporating multi-scale feature extraction, enhanced encoder–decoder connections, and refined upsampling to better handle complex, high-resolution imagery.
To obtain larger receptive fields without lowering image resolution via downsampling, DeepLab was introduced in \cite{deeplabv1}, utilizing atrous convolutions to achieve dense, multi-scale feature extraction. The work in \cite{10608051} further improves DeepLab by aggregating multi-scale features to enhance aerial image segmentation. Many early segmentation models did not fully exploit low-level features in the decoder, resulting in the loss of fine spatial detail. U-Net \cite{ronneberger2015u} addresses this limitation by introducing skip connections that link encoder and decoder features, enabling more effective fusion of detailed and high-level information. U-Net and its variants have also been widely adopted for UAV and aerial image segmentation, as in UVid-Net for UAV video segmentation \cite{9392319} and Res-U-Net-based coastal boundary detection from high-resolution UAV imagery \cite{wang2025multi}.

Attention-based architectures have also been used for aerial image segmentation. SegHSI \cite{10751785} adopts a CNN-based backbone augmented with cluster attention to model spectral–spatial correlations in hyperspectral imagery. The work in \cite{10422983} employs a Transformer–CNN hybrid encoder–decoder for UAV forestry segmentation, enabling accurate delineation of diseased pine canopies in high resolution imagery.


\subsection{Rotation invariance in remote sensing imagery} \label{subsec:seg-rsd}
Rotation invariance (or equivariance) is crucial for accurate object-level segmentation in UAV imagery, as targets often appear in arbitrary orientations.
The work in \cite{isprs-annals-V-1-2022-15-2022} introduces a CNN-based framework for vehicle instance segmentation in UAV images that uses rotated bounding boxes instead of standard axis-aligned boxes to better match the true orientation of vehicles.
The method in \cite{9466361} learns rotation-invariant deep embeddings by aligning features from multiple rotated views of the same image, enabling robust analysis of remote sensing images under arbitrary orientations. FRINet, introduced in \cite{10339376}, is built on CNN feature extractors and incorporates a rotation-adaptive matching mechanism to enable few-shot, rotation-invariant segmentation of aerial images. The paper \cite{mo2024achieving} proposes a mechanism-assured convolution that embeds rotation invariance directly into the operation itself, avoiding reliance on data augmentation and ensuring consistent feature responses under rotated inputs.
RIC-CNN \cite{mo2024ric} achieves rotation invariance by embedding polar coordinate information into convolution operations, enabling the network to generate stable feature representations under arbitrary rotations.
The method in \cite{mitton2021rotation} embeds rotation equivariance into UNet by using the discrete rotation group $C_8$ ($45^\circ$ increments) as the symmetry group, enabling stable feature representations that improve both deforestation segmentation and driver classification under arbitrary image orientations.
In this paper, we introduce a computationally efficient rotation-invariant convolution and incorporate it into a U-Net architecture for semantic segmentation of UAV imagery. 

\subsection{Single-orientation rotation convolution}\label{sec:dense-conv} \label{subsec:standard-conv}
As single-orientation convolution reduces to standard convolution, we review efficient implementations of standard convolution in this section. Several FFT-based approaches \cite{FFT1,FFT2,FFT3,FFT4} have been proposed to reduce the computational cost for standard convolution on different hardware. FFT-based approaches can significantly lower the work complexity of convolution, but it is only effective for convolution with larger filters \cite{cuDNN}. 

Lavin \& Gray \cite{Lavin} proposed a method based on the Winograd algorithm \cite{winograd} for convolution with small filters. The key idea of Winograd-based convolution is similar to the one based on FFT which performs multiplications in the complex domain. However, unlike FFT-based convolution, the Winograd-based convolution operates on real numbers, thus requiring fewer operations. 
For a filter with larger size, Winograd-based convolution shows poor performance, since it needs to perform a significant amount of extra arithmetic such as additions and data transformations.

To reduce the substantial memory usage in convolution with matrix multiplication, Lu \textit{et al.} \cite{lu2023im2win} proposed new data-lowering techniques. However, their method still requires data duplication. In \cite{10.5555/3213069.3213072}, Chang \textit{et al.} proposed an approach to improve memory efficiency in convolution operations by effectively utilizing the on-chip memory of GPUs. In \cite{9235051}, a matrix multiplication-based approach without data lowering is proposed. However, it requires recomputing writing positions, preventing direct pointwise operations.


\subsection{Rotation-equivariant convolution for multi-orientation rotations} \label{subsec:rot-conv}
A general theory of an equivariant CNN can be found in \cite{cohen2018general}. The mathematical formulation of equivariance through a weight sharing scheme is exhibited in \cite{equi-theory}.


G-convolution \cite{G} applies a group of filters to each feature map, where each filter is a transformed version of the original filter. Two different groups are proposed in \cite{G}, code-name $p4$ group and $p4m$ group. The $p4$ group consists of 4 filters where each filter is obtained by rotating the original filter by 0 degrees, 90 degrees, 180 degrees and 270 degrees, respectively. The $p4m$ group consists of 8 filters which are obtained by combining $p4$ group and reflection transformations. The convolution with $p4$ group or $p4m$ group will increase the computational overhead by 4 or 8 times. 
There are multiple works \cite{CubeNet1, CubeNet2} applying the G-convolution on 3D volume data. In \cite{CubeNet1}, different groups code-name \textit{Cube Group}, \textit{Klein's Four-group} and \textit{Tetrahedral Group} are proposed. Among these groups, the largest group is the \textit{Cube Group} which consists of 24 filters resulting in an increase of computational overhead by 24 times. Obviously, the number of filters in a group will increase along with increasing dimension of the input data and therefore the computational overhead. 
Hexagonal group convolution \cite{hexa} rotates a filter by any multiple of 60 degrees, without interpolation. They define the following two hexagonal group convolutions: groups $p6$ and $p6m$, which contain integer translations (orientation cycling), rotations with multiples of 60 degrees, and mirroring for $p6m$. The paper \cite{G-cnn10} proposes Gauge-equivariant convolutional network for manifolds, which uses Hexagonal group convolution to implement the weight sharing scheme for Gauge-equivariant. 
The work in \cite{shallowNet}
proposes a shallow network which
performs convolution with rotated versions
of each canonical filter, where the rotation is with arbitrary angle.
The convolution is followed by a global pooling over orientations.
The number of rotations applied is 32 i.e. the angle increment between 
successive two rotations is $\frac{\pi}{16}$.
Bicubic interpolation is used for sampling after rotation.


Oriented Response Network (ORN) \cite{orn} introduces active rotating filters that generate orientation-aware feature maps, allowing CNNs to model rotations explicitly and achieve rotation invariance or equivariance without heavy data augmentation.
The work in \cite{e2cnn} provides a unified framework for designing steerable CNNs that are equivariant to all
$E(2)$ transformations and offers the widely used E2CNN library for generating rotation-equivariant features.

Filters with steerability can be
reconstructed at any orientation as a finite, linear combination of base
filters. This eliminates the need to learn multiple filters at different
orientations.
Harmonic Network proposed in \cite{harmonic-cnn} 
achieves rotation equivariance in continuous rotations by  
learning coefficients of base filters, where
these base filters are used to form the filters (steerable) of the CNN model.
In the network proposed in \cite{G-cnn7},
rotation equivariance is guaranteed by using 
the G-convolution in which the number of rotation is more than four.
The rotated filter (steerable) in a group is constructed by manipulating coefficients of the base filter. 
In this way, it avoids the use of interpolation which usually introduces artifacts in the rotated filter.  
The network shown in \cite{gabor} achieves rotation equivariance by multiplying each convolutional filter by 
several oriented Gabor filters \cite{weldon1996efficient}. In the back propagation stage, only the convolutional filters are updated.
The method in \cite{cheng2018rotdcf} reduces the number of parameters and computational complexity of rotation-equivariant
CNNs by decomposing convolutional filters (steerable) under joint
steerable bases over space and rotations simultaneously. Meanwhile, the performance is preserved. 

Rotation equivariance typically leads to a rapid increase in feature dimensionality as the number of sampled orientations grows. To tackle this problem of exploding dimensionality, a rotation equivariant convolution is usually followed by an orientation pooling operation, after which the layer becomes rotation invariant. The orientation pooling can be implemented in two different ways: point-wise average pooling \cite{featureTrans} or point-wise max pooling \cite{shallowNet}.

In this paper, we describe how our proposed convolution operation is applied to accelerate a rotation invariant convolution layer, which combines symmetric rotation and steerable transformations. 